\documentclass[12pt]{article}
\usepackage[T1]{fontenc}
\usepackage[utf8]{inputenc}
\usepackage{lmodern}
\usepackage{textcomp}
\usepackage{enumitem}
\usepackage{geometry}
\geometry{margin=1in}
\begin{document}
\begin{center}
Answer Key - Business Administration - Undergraduate Level Assessment 3 \
\small (Answers are ordered exactly as in the question paper.)
\end{center}

\section*{Multiple Choice}
\begin{enumerate}[label=\arabic*.]
\item \textbf{Answer:} C
\\ \textit{Options:} A) Give the reporter other information she is certain is correct.; B) Say that the information is 'off the record' and will be disseminated later.; C) Say 'I don't know' and promise to provide the information later.; D) Say 'no comment,' rather than appear uninformed.
\\ \textit{Note:} This response demonstrates professionalism and honesty, acknowledging the limitation of knowledge while committing to follow up with accurate information, which is essential for maintaining credibility in public relations.
\item \textbf{Answer:} B
\\ \textit{Options:} A) Legal, Interactions, Mechanics; B) Civil, Relations, Outcomes; C) Political, Interactions, Outcomes; D) Environmental, Relations, Mechanics
\\ \textit{Note:} The correct answer is "Civil, Relations, Outcomes" because civil regulation encompasses not only the relationships that civil society organizations (CSOs) have with businesses but also addresses the outcomes and impacts of these interactions on broader societal issues.
\item \textbf{Answer:} D
\\ \textit{Options:} A) Functional; B) Operational; C) Middle level; D) Top level
\\ \textit{Note:} A corporate manager operates at the top level of an organization because they are responsible for setting strategic direction, making high-level decisions, and overseeing the overall management of the organization.
\item \textbf{Answer:} C
\\ \textit{Options:} A) Sponsorships.; B) Exhibitions.; C) Branded content.; D) Advertisement.
\\ \textit{Note:} Branded content is the correct answer because it specifically refers to entertainment material created by a company or brand that is distributed through paid or owned media channels to promote itself.
\item \textbf{Answer:} B
\\ \textit{Options:} A) Work flow design; B) Work schedule design; C) Work rate design; D) Work output design
\\ \textit{Note:} A Gantt chart is a type of work schedule design because it visually represents project timelines, tasks, and their durations, allowing for effective planning and tracking of work progress.
\end{enumerate}

\section*{True / False}
\begin{enumerate}[label=\arabic*.]
\setcounter{enumi}{5}
\item \textbf{Answer:} False
\\ \textit{Options:} A) True; B) False
\\ \textit{Note:} Cost-based pricing focuses on the costs of production and adding a markup, rather than starting with the customer's perception of value.
\item \textbf{Answer:} True
\\ \textit{Options:} A) True; B) False
\\ \textit{Note:} Community relations involve the development and implementation of strategies that enhance collaboration, communication, and mutual understanding between an organization and its local community, thereby fostering positive relationships.
\item \textbf{Answer:} True
\\ \textit{Options:} A) True; B) False
\\ \textit{Note:} Donaldson and Preston (1995) categorize stakeholder theory into three distinct types-Normative, Descriptive, and Instrumental-each serving different purposes in understanding stakeholder relationships and their implications for business practices.
\item \textbf{Answer:} False
\\ \textit{Options:} A) True; B) False
\\ \textit{Note:} The statement is false because an objective must be specific, measurable, achievable, relevant, and time-bound (SMART) to effectively guide actions and assess progress.
\item \textbf{Answer:} True
\\ \textit{Options:} A) True; B) False
\\ \textit{Note:} Programmed decisions are considered true as they involve repetitive situations where established procedures and guidelines can be consistently applied to achieve standard outcomes.
\end{enumerate}

\section*{Long Form Response}
\begin{enumerate}[label=\arabic*.]
\setcounter{enumi}{10}
\item \textbf{Answer:} Motive development in consumer behavior refers to the process through which consumers identify their needs and desires, leading to the recognition of a problem that requires resolution. This concept is rooted in the imbalance between a consumer's actual state-what they currently possess or experience-and their desired state-what they aspire to achieve or acquire. When consumers perceive a gap between these two states, it generates a motivational drive to address the discrepancy, prompting them to seek solutions through purchasing decisions. This recognition of a problem often triggers an evaluation of options, as consumers aim to restore balance and satisfy their unmet needs. Ultimately, understanding motive development is crucial for marketers, as it informs strategies to effectively engage consumers and influence their buying behavior.
\\ \textit{Note:} correct because it effectively outlines how the disparity between a consumer's current situation and their goals creates a motivational drive that leads to problem recognition, which is central to understanding motive development in consumer behavior.
\item \textbf{Answer:} Transformational leaders possess a strategic thinking and outlook that emphasizes long- term vision and adaptability to change, which significantly influences organizational change and employee motivation. By articulating a clear and inspiring vision, these leaders align their team's efforts with organizational goals, fostering a sense of purpose among employees. Their ability to anticipate challenges and innovate solutions encourages a culture of continuous improvement and engagement. Furthermore, transformational leaders empower employees by involving them in decision-making processes, which enhances motivation and commitment to the organization's success. Overall, their strategic approach not only drives effective change but also cultivates a motivated workforce ready to embrace transformation.
\\ \textit{Note:} correct because it highlights how transformational leaders foster a long-term vision and adaptability, which not only aligns organizational goals with employee efforts but also cultivates a culture of engagement and continuous improvement, ultimately driving change and motivation within the organization.
\end{enumerate}

\vfill
\noindent\textit{End of Answer Key}
\end{document}